% Digitalización e impuestos: ¿qué cambia cuando el gobierno sabe más? (Jacobs, 2017)
% Finanzas Públicas, Maestría en Economía Aplicada, Tec de Monterrey
% Esquema aprobado: esquema_jacobs.md (2026-10-26). Compilar con: latexmk -pdf jacobs_digitalizacion.tex
% Versión para repartir: jacobs_digitalizacion_handout.tex (misma fuente, opción handout).
% Las notas del presentador van en \note{}; para verlas, compilar con \setbeameroption{show notes}.
\documentclass[aspectratio=169,11pt]{beamer}

\usepackage[T1]{fontenc}
\usepackage{lmodern}
\usepackage[spanish,es-nodecimaldot]{babel}
\unaccentedoperators % operadores sin acento (la opción es-noaccentedoperators no existe en esta versión de babel)
\usepackage{amsmath,amssymb}
\usepackage{booktabs}
\usepackage{array}
\usepackage{tikz}
\usetikzlibrary{babel,arrows.meta,calc,positioning,decorations.pathreplacing,patterns}
\usepackage{pgfplots}
\pgfplotsset{compat=1.18}

% ---------------------------------------------------------------------------
% Tema sobrio (igual al de las sesiones anteriores)
% ---------------------------------------------------------------------------
\definecolor{cTitulo}{RGB}{0,55,100}
\definecolor{cOferta}{RGB}{31,87,153}
\definecolor{cDemanda}{RGB}{200,100,0}
\definecolor{cImpuesto}{RGB}{175,30,30}
\definecolor{cRecauda}{RGB}{40,130,80}
\definecolor{cGris}{RGB}{110,110,110}

\usetheme{default}
\setbeamertemplate{navigation symbols}{}
\setbeamercolor{frametitle}{fg=cTitulo}
\setbeamercolor{title}{fg=cTitulo}
\setbeamercolor{structure}{fg=cTitulo}
\setbeamercolor{block title}{bg=cTitulo!12,fg=cTitulo}
\setbeamercolor{block body}{bg=cTitulo!4}
\setbeamercolor{block title alerted}{bg=cImpuesto!12,fg=cImpuesto}
\setbeamercolor{block body alerted}{bg=cImpuesto!4}
\setbeamerfont{frametitle}{size=\large,series=\bfseries}
\setbeamertemplate{itemize items}[circle]
\setbeamertemplate{blocks}[rounded]
\setbeamertemplate{footline}{%
  \hfill{\small\color{cGris}Digitalización · Jacobs (2017)\quad\insertframenumber/\inserttotalframenumber}\hspace*{1em}\vspace*{0.6em}}

% Fuente al pie de cada diapositiva
\newcommand{\fuente}[1]{\par\vfill{\small\color{cGris}#1\par}}
% Resultado final de una derivación, resaltado
\newcommand{\resultado}[1]{\colorbox{cTitulo!12}{$\displaystyle #1$}}

% ---------------------------------------------------------------------------
% Estilo común de gráficas
% ---------------------------------------------------------------------------
\tikzset{
  every node/.append style={font=\small},
  eje/.style={-{Stealth[length=2mm]},thick,black!75},
  flecha/.style={-{Stealth[length=2mm]},thick,cTitulo},
  guia/.style={thin,dashed,black!45},
  punto/.style={circle,fill=black,inner sep=1.2pt},
  caja/.style={draw=cTitulo,rounded corners,fill=cTitulo!5,align=center,inner sep=4pt},
  cajafisco/.style={draw=cRecauda,rounded corners,fill=cRecauda!10,align=center,inner sep=3pt},
}

\title{Digitalización e impuestos: ¿qué cambia cuando el gobierno sabe más?}
\subtitle{Digitalización de las finanzas públicas · Segunda sección}
\author{Héctor Juan Villarreal}
\institute{Finanzas Públicas · Maestría en Economía Aplicada\\Escuela de Gobierno y Transformación Pública · Tec de Monterrey}
\date{Octubre de 2026}

\begin{document}

% =========================================================================
% I. EL PROBLEMA DE FONDO ES DE INFORMACIÓN · 22 min
% =========================================================================
\section{Información}

% 1 -----------------------------------------------------------------------
\begin{frame}[plain]
  \titlepage
  \note{Lectura: Jacobs, B. (2017), «Digitalization and Taxation», cap.~2 en S. Gupta, M. Keen, A. Shah y G. Verdier (eds.), \emph{Digital Revolutions in Public Finance}, FMI, pp.~25--55. El capítulo es técnico; hoy el peso está en las aplicaciones.}
\end{frame}

% 2 -----------------------------------------------------------------------
\begin{frame}{De Kaplow a Jacobs: el problema es de información}
  \begin{columns}[T]
    \begin{column}{0.48\textwidth}
      \begin{block}{Lo que nos dejó Kaplow}
        \begin{itemize}
          \item En Mirrlees, el gobierno ve el ingreso, no la capacidad (Kaplow, diap.~27).
          \item En México, el instrumento para redistribuir es débil (Kaplow, diap.~17).
        \end{itemize}
      \end{block}
    \end{column}
    \begin{column}{0.48\textwidth}
      \pause
      \begin{block}{Lo que propone Jacobs}
        La digitalización puede aflojar las restricciones de información. Nunca las elimina.
      \end{block}
    \end{column}
  \end{columns}
  \pause
  \vspace{1em}
  \begin{alertblock}{Pregunta para hoy}
    ¿Qué cambia cuando el gobierno sabe más?
  \end{alertblock}
  \fuente{Jacobs, introducción, pp.~25--26.}
\end{frame}

% 3 -----------------------------------------------------------------------
\begin{frame}{Con información perfecta, no habría disyuntiva}
  \begin{itemize}
    \item Si el gobierno viera todo sin costo, cobraría \textbf{impuestos de suma fija individualizados}: según la capacidad, las necesidades, la herencia o la suerte de cada quien.
    \item Tampoco habría evasión: sabría cuánto gana, ahorra y consume cada persona.
  \end{itemize}
  \pause
  \begin{block}{Segundo teorema del bienestar}
    Con mercados perfectos, cualquier asignación eficiente se alcanza con mercados y transferencias de suma fija. La eficiencia y la distribución se deciden por separado.
  \end{block}
  {\small\color{cGris}Repaso: el primer teorema dice que el equilibrio competitivo es eficiente. Jacobs solo usa el segundo.}
  \pause
  \begin{block}{Enlace}
    Kaplow, diap.~6: si todos fuéramos iguales, bastaría la suma fija. Aquí las personas sí difieren, pero el gobierno lo vería todo.
  \end{block}
  \fuente{Jacobs, introducción, p.~25.}
\end{frame}

% 4 -----------------------------------------------------------------------
\begin{frame}{Dos cosas que el gobierno no ve}
  \begin{columns}[T]
    \begin{column}{0.44\textwidth}
      \centering
      \begin{tikzpicture}[scale=0.8]
        \draw[eje] (0,0) -- (5.2,0);
        \draw[eje] (0,0) -- (0,4.8);
        \node[below] at (2.6,0) {bienestar del menos capaz};
        \node[above,rotate=90] at (0,2.4) {bienestar del más capaz};
        \draw[very thick,cOferta] plot[domain=6:84,samples=60] ({4.5*cos(\x)},{4.3*sin(\x)});
        \draw[very thick,cImpuesto] (1.54,4.04) .. controls (2.4,3.6) and (2.95,2.7) .. (3.0,1.1);
        \onslide<3->{\draw[very thick,cRecauda,dashed] (1.54,4.04) .. controls (2.9,3.75) and (3.6,2.9) .. (3.75,1.6);}
        \node[punto] at (1.54,4.04) {};
        \node[anchor=south west] at (1.3,4.12) {sin redistribuir};
        \draw[very thick,cOferta] (0.2,-0.95) -- +(0.5,0) node[right,black] {primer mejor};
        \draw[very thick,cImpuesto] (0.2,-1.4) -- +(0.5,0) node[right,black] {segundo mejor};
        \onslide<3->{\draw[very thick,cRecauda,dashed] (0.2,-1.85) -- +(0.5,0) node[right,black] {con más información};}
      \end{tikzpicture}
    \end{column}
    \begin{column}{0.54\textwidth}
      \begin{block}{(a) El resultado verdadero}
        Cuánto ganó o consumió cada quien: de ahí la evasión. Como las multas no pueden ser infinitas, manda la \textbf{tecnología de fiscalización}.
      \end{block}
      \pause
      \begin{block}{(b) Las características y la conducta}
        La capacidad y el esfuerzo: de ahí la disyuntiva entre equidad y eficiencia de Mirrlees.
      \end{block}
      \onslide<3->{Más información empuja hacia afuera la frontera de segundo mejor. No la lleva a la de primer mejor.}
    \end{column}
  \end{columns}
  \fuente{Jacobs, pp.~25--28 y 32--33; Allingham y Sandmo (1972) y Mirrlees (1999), citados por Jacobs. Gráfica propia.}
  \note{Las dos fronteras se tocan en el punto sin redistribución: si el gobierno no quiere redistribuir, no necesita saber la capacidad de nadie.}
\end{frame}

% 5 -----------------------------------------------------------------------
\begin{frame}{En el modelo de Mirrlees, la digitalización no sirve de nada}
  \begin{itemize}
    \item Mirrlees supone que el ingreso se ve sin error y que la capacidad no se ve. Más datos no cambian ninguna de las dos cosas.
    \item Con los supuestos de Atkinson-Stiglitz pasa lo mismo: todo se redistribuye con el ISR (Kaplow, diap.~13).
    \item La promesa de la digitalización depende de que esos supuestos fallen. Y en la realidad fallan.
  \end{itemize}
  \pause
  \begin{block}{¿Y si algún día se pudiera medir la capacidad?}
    Jacobs duda, por cuatro razones. No está claro qué medir: el salario por hora ya depende de la educación y del esfuerzo. Hasta el IQ cambia. Medirla abre problemas filosóficos, políticos y legales. Y quien sepa que lo van a gravar por su capacidad, la va a esconder.
  \end{block}
  \fuente{Jacobs, pp.~33--34 (nota~9) y 36--37.}
\end{frame}

% 6 -----------------------------------------------------------------------
\begin{frame}{Dos canales: fiscalizar mejor y diseñar mejor}
  \centering
  \begin{tikzpicture}[node distance=5mm and 9mm]
    \node[caja,text width=3.0cm] (d) {Digitalización:\\más información};
    \node[caja,text width=5.6cm,right=of d,yshift=9mm] (c1) {\textbf{Canal 1. Fiscalizar mejor}\\ver lo que ya se grava (5 ideas)};
    \node[caja,text width=5.6cm,right=of d,yshift=-9mm] (c2) {\textbf{Canal 2. Diseñar mejor}\\condicionar el impuesto en más información (8 ideas)};
    \draw[flecha] (d.east) -- (c1.west);
    \draw[flecha] (d.east) -- (c2.west);
    \node[right=3mm of c1,text width=2.0cm,align=left,color=cGris] {diaps.~7--11};
    \node[right=3mm of c2,text width=2.0cm,align=left,color=cGris] {diaps.~12--16};
  \end{tikzpicture}
  \par\medskip\raggedright
  \pause
  En los dos canales: la misma recaudación con tasas más bajas, o más redistribución con las mismas tasas.
  \pause
  \begin{block}{Ojo}
    La digitalización privada también ayuda a eludir, con planeación fiscal más agresiva. Por eso no es obvio si las tasas óptimas suben o bajan.
  \end{block}
  \fuente{Jacobs, pp.~26--27 y conclusión, pp.~50--51. Diagrama propio.}
\end{frame}

% =========================================================================
% II. CANAL 1: FISCALIZAR MEJOR · 15 min
% =========================================================================
\section{Fiscalizar mejor}

% 7 -----------------------------------------------------------------------
\begin{frame}{La evasión cae cuando alguien más informa}
  \renewcommand{\arraystretch}{1.25}
  \begin{tabular}{@{}>{\raggedright\arraybackslash}p{0.34\textwidth}>{\raggedright\arraybackslash}p{0.56\textwidth}@{}}
    \toprule
    \textbf{Quién informa} & \textbf{Qué ve el fisco} \\
    \midrule
    El patrón & El salario de sus empleados \\
    Bancos, fondos, aseguradoras & Intereses, saldos y rendimientos \\
    El pago con tarjeta & Las ventas de la empresa: el consumidor reporta \\
    Nadie & El ingreso del trabajador por cuenta propia \\
    \bottomrule
  \end{tabular}
  \par\medskip
  \pause
  \begin{itemize}
    \item Quien evade prefiere el efectivo. Si mucho se paga en efectivo, la digitalización ayuda poco.
    \item Sin efectivo, el trueque se vuelve negocio.
  \end{itemize}
  \pause
  Enlace: la informalidad vuelve más elástica la base formal (26 de septiembre, diap.~14).
  \fuente{Jacobs, «Digitalization and Tax Enforcement», pp.~28 y 31--32; nota~3 (p.~29).}
\end{frame}

% 8 -----------------------------------------------------------------------
\begin{frame}{México: cada factura pasa por el SAT}
  \centering
  \begin{tikzpicture}[x=1cm,y=1cm]
    \node[caja,text width=2.3cm] (v) at (0,0) {Vendedor\\emite el CFDI};
    \node[cajafisco,text width=4.0cm] (s) at (4.6,0) {SAT o proveedor\\autorizado: valida\\y asigna folio};
    \node[caja,text width=2.3cm] (c) at (9.2,0) {Comprador\\lo usa para deducir};
    \draw[flecha] (v) -- (s);
    \draw[flecha] (s) -- (c);
    \onslide<2->{
      \node[cajafisco,text width=6.6cm] (x) at (4.6,-1.45) {El SAT cruza las dos puntas: lo que uno declara vender, el otro lo declara comprar};
      \draw[flecha,cRecauda] (v.south) |- (x.west);
      \draw[flecha,cRecauda] (c.south) |- (x.east);}
  \end{tikzpicture}
  \par\medskip\raggedright
  \begin{itemize}
    \item Desde el 1 de enero de 2014, todo comprobante fiscal es una factura electrónica: el CFDI.
    \item Antes de expedirla, se manda al SAT o a un proveedor que él autoriza. El comprador está obligado a pedirla.
    \onslide<3->{\item Lectura con el cobro fraccionado: cada factura deja huella (sesión del IVA, diap.~5).}
  \end{itemize}
  \fuente{CFF, art.~29 (texto vigente); decreto de reformas al CFF, DOF 9-12-2013, transitorio primero. Respaldo conceptual: Jacobs, p.~32.}
\end{frame}

% 9 -----------------------------------------------------------------------
\begin{frame}{Brasil: el impuesto se cobra al momento de pagar}
  \begin{columns}[T]
    \begin{column}{0.56\textwidth}
      \centering
      \begin{tikzpicture}[x=1cm,y=1cm]
        \node[caja,text width=1.7cm] (comp) at (0,0) {Comprador\\paga 126.5};
        \node[caja,text width=1.6cm] (ent) at (2.9,0) {Entidad\\de pago};
        \node[cajafisco,text width=1.6cm] (fis) at (2.9,-2.2) {Fisco};
        \node[caja,text width=1.6cm] (prov) at (5.6,0) {Proveedor};
        \draw[flecha] (comp) -- (ent);
        \draw[flecha,cRecauda,{Stealth[length=2mm]}-{Stealth[length=2mm]}] ([xshift=-4mm]ent.south) -- node[left,align=right] {¿tiene\\créditos?} ([xshift=-4mm]fis.north);
        \onslide<2->{
          \draw[flecha,cRecauda] ([xshift=4mm]ent.south) -- node[right,align=left] {si no:\\26.5} ([xshift=4mm]fis.north);
          \draw[flecha] (ent) -- node[above] {100} (prov);}
      \end{tikzpicture}
    \end{column}
    \begin{column}{0.42\textwidth}
      \begin{itemize}
        \item La factura dice cuánto se debe; el pago fraccionado lo cobra en el momento.
        \onslide<2->{\item Si el proveedor no tiene créditos, la entidad de pago separa los 26.5 y le pasa 100. Si los tiene, le pasa todo.}
        \onslide<3->{\item El cumplimiento se vuelve casi automático. Brasil ya tenía factura electrónica y pagos instantáneos (Pix).}
      \end{itemize}
    \end{column}
  \end{columns}
  \par\medskip
  {\small El 26.5\,\% es una tasa estimada, no fijada en ley. Ya lo vimos en la sesión del IVA, diap.~20.}
  \fuente{Barreix, Rocha y Tostes Neto (2025), nota larga, gráfico~3 (redibujado); LC 214/2025, arts.~31--35. Respaldo conceptual: Jacobs, p.~32.}
\end{frame}

% 10 ----------------------------------------------------------------------
\begin{frame}{El presupuesto de cada persona delata lo que no declara}
  \[
    \underbrace{\Delta a}_{\text{ahorro}} + \underbrace{b^*}_{\text{herencias recibidas}} = \underbrace{r^*a}_{\text{capital}} + \underbrace{w^*l}_{\text{trabajo}} - \underbrace{p^*c}_{\text{consumo}} - \underbrace{g}_{\text{regalos hechos}}
  \]
  {\small\color{cGris}El asterisco marca valores después de impuestos. Jacobs usa aquí $g$ para los regalos; en la diap.~12, para una transferencia.}
  \begin{itemize}
    \item Si el gobierno ve el consumo y la riqueza, puede revisar si el ingreso declarado alcanza.
    \item A lo largo de la vida, $C + B = Y$: si el consumo se aleja del ingreso, hay herencias o regalos sin declarar.
  \end{itemize}
  \pause
  \begin{block}{México: la discrepancia fiscal}
    Si una persona gasta, compra o deposita en un año más de lo que declaró, el SAT presume que la diferencia es ingreso omitido. Aplica también a quien no está inscrito en el RFC.
  \end{block}
  \fuente{Jacobs, pp.~28--30; LISR, art.~91 (texto vigente).}
\end{frame}

% 11 ----------------------------------------------------------------------
\begin{frame}{El capital ya no se esconde tan fácil en el extranjero}
  \begin{itemize}
    \item Con el intercambio automático de información, cada fisco recibe de los otros países los saldos y los intereses que sus residentes tienen en bancos de allá.
    \item México se comprometió a intercambiar desde 2017. En 2024 mandó datos a 79 jurisdicciones.
  \end{itemize}
  \pause
  \begin{itemize}
    \item Con registros que se cruzan, el capital se grava donde vive su dueño. El ISR de las empresas podría quedar como una retención sobre dividendos.
    \item Jacobs propone un solo impuesto a todo el ingreso del capital, con una tasa y una exención: se acaba el arbitraje entre activos.
  \end{itemize}
  \pause
  \begin{block}{Límite}
    Beneficiarios finales ocultos y títulos al portador.
  \end{block}
  \fuente{Jacobs, pp.~31--32 y 39--43; OCDE (2024), \emph{Peer Review of the Automatic Exchange of Financial Account Information 2024 Update}, cuadro~1.1.}
\end{frame}

% =========================================================================
% III. CANAL 2: DISEÑAR MEJOR · 17 min
% =========================================================================
\section{Diseñar mejor}

% 12 ----------------------------------------------------------------------
\begin{frame}{Un IVA progresivo sin impuesto al ingreso}
  El gobierno no ve el ingreso, solo grava el consumo con una tasa única $\tau$.
  \begin{align*}
    \onslide<1->{\text{Sin transferencia:}\quad & wl = (1+\tau)\,c \quad\Rightarrow\quad \text{paga } \tau c\text{, proporcional}\\}
    \onslide<2->{\text{Con transferencia uniforme } g\text{:}\quad & wl + g = (1+\tau)\,c\\}
    \onslide<3->{\text{Pago neto:}\quad & \tau c - g\\}
    \onslide<4->{\text{Tasa media neta:}\quad & \resultado{\tau - \frac{g}{c}} \quad \text{sube con el consumo}}
  \end{align*}
  \onslide<5->{
  \begin{block}{Resultado}
    Con identificación biométrica y pagos electrónicos, una transferencia igual para todos vuelve progresivo un impuesto proporcional al consumo. Y le quita la razón de ser a la tasa cero en alimentos (sesión del IVA, diap.~16).
  \end{block}
  {\small Es la cuenta del impuesto plano de Kaplow, diap.~5. Notación: $w$, $l$, $c$ y $g$ como en Kaplow; aquí $\tau$ es una tasa única al consumo, no el vector $\tau_i$.}}
  \fuente{Jacobs, pp.~34--35, notas~10 y~11.}
\end{frame}

% 13 ----------------------------------------------------------------------
\begin{frame}{El IVA personalizado: devolver solo a quien lo necesita}
  \begin{columns}[T]
    \begin{column}{0.52\textwidth}
      \begin{itemize}
        \item La $g$ de Jacobs es igual para todos. Barreix la focaliza: devuelve a los hogares de menor ingreso el IVA de una canasta típica (sesión del IVA, diap.~19).
        \item Brasil: cashback para familias con ingreso per cápita de hasta medio salario mínimo. Antes, el Devolve ICMS de Rio Grande do Sul: desde 2021, más de 600 mil familias.
      \end{itemize}
    \end{column}
    \begin{column}{0.45\textwidth}
      \begin{block}{Lo que dice el FMI}
        Recuadro 1: un impuesto progresivo al consumo necesita saber quién compra, y el IVA es anónimo.\\[3pt]
        Nota 28: el IVA personalizado de Barreix y el IVA progresivo de de la Feria y Swistak, que devuelve en el momento de la compra.
      \end{block}
    \end{column}
  \end{columns}
  \pause
  \par\smallskip
  Es una respuesta a la pregunta de Kaplow, diap.~17.
  \pause
  \begin{alertblock}{Pregunta}
    ¿A quién le devolverían, si buena parte de los hogares vive de la informalidad?
  \end{alertblock}
  \fuente{de Mooij, Hebous y Keen (2025), IMF WP/25/165, recuadro~1 y nota~28; Barreix, Rocha y Tostes Neto (2025), nota larga; LC 214/2025, arts.~113 y 118. Jacobs, pp.~34 y 51.}
  \note{Cashback brasileño: 100\,\% de la CBS y 20\,\% del IBS en gas, luz, agua, drenaje y telecomunicaciones; 20\,\% en lo demás (sesión del IVA, diap.~20). El Devolve ICMS paga devoluciones trimestrales.}
\end{frame}

% 14 ----------------------------------------------------------------------
\begin{frame}{Donde no hay reventa, el precio puede depender de cuánto se consume}
  \begin{columns}[T]
    \begin{column}{0.42\textwidth}
      \centering
      \begin{tikzpicture}[x=1cm,y=1cm]
        \fill[cOferta!55] (0.05,0) rectangle (1.5,0.855);
        \fill[cOferta!70] (1.6,0) rectangle (3.05,1.034);
        \fill[cImpuesto!60] (3.15,0) rectangle (4.6,3.026);
        \node[above] at (0.78,0.855) {0.86};
        \node[above] at (2.33,1.034) {1.03};
        \node[above] at (3.88,3.026) {3.03};
        \node[below] at (0.78,0) {básico};
        \node[below] at (2.33,0) {intermedio};
        \node[below] at (3.88,0) {excedente};
        \node at (2.33,-0.75) {más kWh al mes $\to$};
        \draw[eje] (0,0) -- (5.0,0);
        \draw[eje] (0,0) -- (0,3.7) node[above right] {pesos por kWh};
      \end{tikzpicture}
      \par{\small Tarifa doméstica 1, enero de 2021}
    \end{column}
    \begin{column}{0.56\textwidth}
      \begin{itemize}
        \item Un impuesto que cambia según cuánto compra cada quien solo funciona si no hay reventa: el que paga menos le revendería al que paga más.
        \item Agua, luz y gas no se transportan ni se guardan. Se pueden cobrar por bloques.
        \pause
        \item En México, el kWh excedente de la tarifa 1 costaba 3.5 veces el básico.
      \end{itemize}
      \pause
      \begin{block}{Lectura}
        Si el gobierno ve el consumo de cada hogar, puede gravar estos bienes de forma no lineal y redistribuir con menos costo.
      \end{block}
    \end{column}
  \end{columns}
  \fuente{Jacobs, «Digitalization and Commodity Taxation», pp.~37--38 y nota~15; SHCP, oficio 349-B-1-I-244 (15-12-2020), anexo, tarifa~1.}
\end{frame}

% 15 ----------------------------------------------------------------------
\begin{frame}{Etiquetado: gravar lo que se ve y anda junto con la capacidad}
  \begin{itemize}
    \item El \textbf{etiquetado} (\emph{tagging}, Akerlof, 1978): condicionar el impuesto en rasgos que se verifican, cambian poco y van con la capacidad o la necesidad: edad, sexo, discapacidad, hijos.
    \item Casi todos los países ya lo hacen, con créditos, deducciones y apoyos a ciertos grupos.
    \item Jacobs va más lejos: tarifas de ISR distintas por edad y por sexo. Los estudios que cita encuentran ganancias grandes.
  \end{itemize}
  \pause
  \begin{block}{Los límites}
    Equidad horizontal: dos personas con el mismo ingreso pagarían distinto (Kaplow, diap.~18). Y la privacidad. Por eso casi nadie lo ha hecho.
  \end{block}
  \pause
  \begin{alertblock}{Pregunta}
    ¿Qué etiqueta les parecería legítima y cuál no?
  \end{alertblock}
  \fuente{Jacobs, «Tagging in Nonlinear Tax Schedules», pp.~49--50, y pp.~51--52; Akerlof (1978), citado por Jacobs.}
\end{frame}

% 16 ----------------------------------------------------------------------
\begin{frame}{Mirar la historia, el patrimonio y la pareja}
  \begin{block}{La regla de Jacobs}
    Un sistema óptimo tiene tantas interdependencias como dimensiones en que difieren las personas.
  \end{block}
  \pause
  \begin{columns}[T]
    \begin{column}{0.32\textwidth}
      \textbf{La historia}\\[2pt]
      Gravar el promedio acumulado del ingreso (Vickrey). Quien tiene años buenos y malos deja de pagar de más. Bastan dos declaraciones seguidas.
    \end{column}
    \begin{column}{0.32\textwidth}
      \pause
      \textbf{El patrimonio}\\[2pt]
      Condicionar los apoyos en los activos. Así no conviene ahorrar hoy para cobrar mañana una pensión de invalidez falsa.
    \end{column}
    \begin{column}{0.32\textwidth}
      \pause
      \textbf{La pareja}\\[2pt]
      La tasa de cada uno depende del ingreso del otro. En una simulación: 25\,\% con pareja sin ingreso, 65\,\% si los dos ganan mucho.
    \end{column}
  \end{columns}
  \fuente{Jacobs, «Jointness in Tax Systems» y siguientes, pp.~43--49 (nota~24); Vickrey (1939, 1947), Golosov y Tsyvinski (2006) y Renes y Zoutman (2016), citados por Jacobs.}
  \note{La intuición: si las personas difieren en más de una cosa, pueden responder al impuesto en más de un margen a la vez (trabajar menos y ahorrar más). Amarrar los esquemas entre sí encarece esas jugadas dobles.}
\end{frame}

% =========================================================================
% IV. LÍMITES Y CIERRE · 6 min
% =========================================================================
\section{Cierre}

% 17 ----------------------------------------------------------------------
\begin{frame}{Los riesgos}
  \begin{itemize}
    \item Privacidad, equidad horizontal y abuso del poder del Estado.
    \item Con malas instituciones, más información le sirve a un mal gobierno para lograr mejor sus malos fines.
    \item Un sistema más complejo cuesta más de fiscalizar y se entiende menos.
  \end{itemize}
  \pause
  \begin{block}{México: facturas sin operación detrás}
    Si el SAT detecta que alguien factura sin activos, personal ni capacidad para hacerlo, presume que las operaciones no existen y publica la lista en el DOF y en su página. La factura electrónica no impide la factura falsa: la deja registrada.
  \end{block}
  \pause
  \begin{alertblock}{Pregunta}
    ¿Qué información no le darían al SAT?
  \end{alertblock}
  \fuente{Jacobs, pp.~27, 50 y 51--52; CFF, art.~69-B (texto vigente).}
\end{frame}

% 18 ----------------------------------------------------------------------
\begin{frame}{Lo que la digitalización no resuelve}
  \begin{itemize}
    \item El esfuerzo sigue siendo información privada, aun en un mundo totalmente digital y sin evasión.
    \item La disyuntiva entre equidad y eficiencia se reduce, pero no desaparece.
    \item Donde hay informalidad, pasarse a ella es otra forma de responder al impuesto: la base se vuelve más elástica.
  \end{itemize}
  \pause
  \begin{block}{Lectura}
    La frontera de segundo mejor se mueve hacia afuera. La de primer mejor sigue fuera de alcance (diap.~4).
  \end{block}
  \fuente{Jacobs, pp.~26 y 33--34; nota~7 (p.~32).}
\end{frame}

% 19 ----------------------------------------------------------------------
\begin{frame}{Kaplow y Jacobs: la misma teoría, otro énfasis}
  Coinciden: Mirrlees, bienestarismo, Atkinson-Stiglitz, impuesto personalizado.
  \par\smallskip
  \renewcommand{\arraystretch}{1.15}
  \begin{tabular}{@{}>{\raggedright\arraybackslash}p{0.2\textwidth}>{\raggedright\arraybackslash}p{0.37\textwidth}>{\raggedright\arraybackslash}p{0.37\textwidth}@{}}
    \toprule
    & \textbf{Kaplow} & \textbf{Jacobs} \\
    \midrule
    La pregunta & ¿Cómo combinar lo que hay? & ¿Qué cambia si se sabe más? \\
    Atkinson-Stiglitz & Es la referencia & Los datos lo rechazan \\
    Administración & Una salvedad & Uno de los dos canales \\
    Países en desarrollo & Una nota (Kaplow, diap.~17) & Una propuesta: el IVA progresivo \\
    Instituciones & Fuera del análisis & Cuidado con el mal gobierno \\
    \bottomrule
  \end{tabular}
  \pause
  \begin{alertblock}{Pregunta}
    Si Atkinson-Stiglitz falla en los datos, ¿quién debe probar que diferenciar conviene?
  \end{alertblock}
  \fuente{Kaplow, caps.~1--2, notas~2 y~24; sec.~6.C. Jacobs, pp.~27, 34--38 y 51--52.}
  \note{Kaplow da el método para comparar instrumentos sin que la distribución lo enrede todo; Jacobs explica cómo podría construirse el instrumento personalizado que Kaplow da por hecho. La tensión de fondo es de carga de la prueba: Kaplow trata las fallas de Atkinson-Stiglitz como salvedades a una referencia; Jacobs, como el caso que importa.}
\end{frame}

% 20 ----------------------------------------------------------------------
\begin{frame}{Tres ideas para llevarse}
  \begin{enumerate}
    \item La disyuntiva entre equidad y eficiencia es, en el fondo, un problema de información.
    \pause
    \item La digitalización mueve la frontera de segundo mejor por dos canales: fiscalizar mejor y diseñar mejor. No llega al primer mejor.
    \pause
    \item Los instrumentos nuevos valen lo que valgan las instituciones que los usan.
  \end{enumerate}
  \fuente{Jacobs, conclusión, pp.~50--52.}
\end{frame}

% 21 ----------------------------------------------------------------------
\begin{frame}{Puente: el 3 de noviembre}
  \begin{block}{Cuentas nacionales de transferencia y envejecimiento}
    Jacobs ya pide mirar la edad (impuestos por edad) y la vida completa (el promedio de Vickrey). Las cuentas nacionales de transferencia miden cómo cambian el ingreso y el consumo con la edad, y quién sostiene a quién a lo largo de la vida.
  \end{block}
  \fuente{Jacobs, pp.~45--46 y 49--50. Elaboración propia.}
\end{frame}

\end{document}
